\documentclass[12pt,oneside]{book}

% ============================================================
% METADATOS
% ============================================================
\newcommand{\universidad}{Universidad de Buenos Aires (UBA)}
\newcommand{\facultad}{Facultad de Derecho}
\newcommand{\materia}{Contratos Inmobiliarios}
\newcommand{\titulo}{Compraventa: venta de cosas futuras (venta en pozo) y descripción del inmueble en la escritura}
\newcommand{\autor}{Isaac Edgar Camacho Ocampo}
\newcommand{\anio}{2026}

% ============================================================
% PAQUETES
% ============================================================
\usepackage[utf8]{inputenc}
\usepackage[T1]{fontenc}
\usepackage[provide=*,spanish]{babel}
\usepackage[
    a4paper,
    top=2.5cm,
    bottom=2.5cm,
    left=3cm,
    right=2.5cm,
    headheight=15pt
]{geometry}
\usepackage{amsmath}
\usepackage{amssymb}
\usepackage{mathtools}
\usepackage{graphicx}
\usepackage{float}
\usepackage{caption}
\usepackage{subcaption}
\usepackage{booktabs}
\usepackage{array}
\usepackage{longtable}
\usepackage{tabularx}
\usepackage{enumitem}
\usepackage{xcolor}
\usepackage[most]{tcolorbox}
\usepackage{fancyhdr}
\usepackage{ragged2e}
\usepackage[
    colorlinks=true,
    linkcolor=blue!50!black,
    citecolor=green!40!black,
    urlcolor=blue!60!black,
    pdftitle={\titulo},
    pdfauthor={\autor},
    pdfsubject={Apuntes universitarios},
    bookmarks=true
]{hyperref}

% ============================================================
% CONFIGURACIÓN
% ============================================================
\graphicspath{
    {./img/}
    {./graficos/}
    {./figuras/}
}
\setcounter{tocdepth}{2}
\setcounter{secnumdepth}{0}
\setlength{\parindent}{0pt}
\setlength{\parskip}{0.5em}
\setlist[itemize]{
    topsep=0.3em,
    itemsep=0.2em
}
\setlist[enumerate]{
    topsep=0.3em,
    itemsep=0.2em
}
\clubpenalty=10000
\widowpenalty=10000
\addto\captionsspanish{\renewcommand{\chaptername}{Bloque}}

% ============================================================
% COMANDOS
% ============================================================
\newcommand{\abs}[1]{\left|#1\right|}
\newcommand{\norm}[1]{\left\lVert#1\right\rVert}
\newcommand{\vect}[1]{\mathbf{#1}}
\newcommand{\deriv}[2]{\frac{d#1}{d#2}}
\newcommand{\pderiv}[2]{\frac{\partial#1}{\partial#2}}
\newcommand{\mcuad}{$\mathrm{m}^2$}

% ============================================================
% ENTORNOS / CAJAS
% ============================================================
\newtcolorbox{definition}[1]{
    enhanced,
    breakable,
    title={Definición: #1},
    fonttitle=\bfseries,
    colback=blue!3,
    colframe=blue!50!black,
    boxrule=0.8pt,
    arc=2mm
}

\newtcolorbox{important}{
    enhanced,
    breakable,
    title={Importante},
    fonttitle=\bfseries,
    colback=yellow!8,
    colframe=orange!70!black,
    boxrule=0.8pt,
    arc=2mm
}

\newtcolorbox{example}[1]{
    enhanced,
    breakable,
    title={Ejemplo: #1},
    fonttitle=\bfseries,
    colback=green!4,
    colframe=green!40!black,
    boxrule=0.8pt,
    arc=2mm
}

\newtcolorbox{formula}[1]{
    enhanced,
    breakable,
    title={Fórmula / Regla: #1},
    fonttitle=\bfseries,
    colback=purple!4,
    colframe=purple!50!black,
    boxrule=0.8pt,
    arc=2mm
}

\newtcolorbox{exercise}[1]{
    enhanced,
    breakable,
    title={Ejercicio: #1},
    fonttitle=\bfseries,
    colback=gray!5,
    colframe=gray!60!black,
    boxrule=0.8pt,
    arc=2mm
}

\newtcolorbox{note}{
    enhanced,
    breakable,
    title={Nota},
    fonttitle=\bfseries,
    colback=white,
    colframe=gray!50,
    boxrule=0.6pt,
    arc=2mm
}

% ============================================================
% CONFIGURACIÓN FINAL
% ============================================================
\pagestyle{fancy}
\fancyhf{}
\fancyhead[L]{\nouppercase{\leftmark}}
\fancyhead[R]{\materia}
\fancyfoot[C]{\thepage}
\renewcommand{\headrulewidth}{0.4pt}
\renewcommand{\footrulewidth}{0pt}

\fancypagestyle{plain}{
    \fancyhf{}
    \fancyfoot[C]{\thepage}
    \renewcommand{\headrulewidth}{0pt}
}

% ============================================================
\begin{document}

% ============================================================
% PORTADA
% ============================================================
\begin{titlepage}

\thispagestyle{empty}

\begin{center}

\vspace*{1cm}

{\large \textsc{\universidad}}

\vspace{0.2cm}

{\large \textsc{\facultad}}

\vspace{3cm}

{\Huge \textbf{\materia}}

\vspace{1cm}

\parbox{0.85\textwidth}{\centering {\Large \titulo}}

\vspace{0.8cm}

\textit{Comprender cada concepto es el primer paso para dominarlo.}

\vspace{1.2cm}

\rule{0.70\textwidth}{0.5pt}

\vspace{0.5cm}

{\large Apuntes de estudio}

\vspace{0.5cm}

\rule{0.70\textwidth}{0.5pt}

\vfill

{\large \autor}

\vspace{0.3cm}

{\large \anio}

\end{center}

\vfill

\begin{flushleft}
\textit{Este es un modesto aporte a los estudiantes de ``Contratos Inmobiliarios'' no pretende sustituir el material oficial de la materia.}
\end{flushleft}

\end{titlepage}

% ============================================================
% ÍNDICE
% ============================================================
\tableofcontents

% ============================================================
% CONTENIDO
% ============================================================

\chapter*{Introducción}
\addcontentsline{toc}{chapter}{Introducción}
\markboth{Introducción}{Introducción}

La clase cierra el estudio de la \textbf{cosa} en la compraventa y recorre cuatro ideas:

\begin{itemize}
    \item \textbf{Venta de cosas futuras y venta en pozo:} cómo se vende un departamento que todavía no existe, por qué conviene a comprador y constructor, y qué dice el art.~1131 del CCyC (condición suspensiva o asunción del riesgo).
    \item \textbf{Naturaleza jurídica de la venta en pozo:} las tres posturas (compraventa, contrato de obra, contrato innominado y complejo) y la tensión entre dar y hacer.
    \item \textbf{Descripción de la unidad funcional} en la cláusula primera de la escritura, con su orden lógico y la regla de números y letras (CCyC art.~303; Ley 404 art.~61).
    \item \textbf{Descripción del lote:} medidas, linderos, orientación y nomenclatura catastral.
\end{itemize}

\section{Relevancia profesional del tema}

Quien ejerza como \textbf{abogado o escribano} va a asesorar a compradores y desarrolladores en operaciones <<en pozo>>, que son frecuentes y de encuadre discutido: saber qué contrato se firma, qué riesgo asume cada parte y qué cláusulas proteger define la calidad del asesoramiento. Además, la \textbf{cláusula primera de la escritura} (identificar el objeto) es el trabajo cotidiano de una escribanía y es materia de examen: un error en un elemento esencial no es un detalle menor, sino un \textbf{error fuerte}.

Comprar una vivienda es probablemente la decisión económica más grande de una persona, y muchas se hacen sobre planos: este contenido enseña a distinguir qué se está firmando, por qué el precio es menor, quién soporta el riesgo si la obra no se termina y qué debe decir la escritura sobre lo que se compra. Aunque no se trabaje en el área, permite \textbf{leer un boleto o una escritura sin depender de intermediarios} y hacer preguntas que protegen el patrimonio propio y el de la familia.

\section{Cómo se relacionan los temas}

\begin{example}{El cartel en el terreno baldío}
Imaginar que se pasa frente a un terreno baldío con un cartel gigante: \textit{<<Próximamente: unidades de dos ambientes con balcón>>}, y se decide comprar una unidad. Los cuatro bloques siguen ese camino: \textbf{(1)} entender por qué se vende algo que todavía no existe y qué dice el Código; \textbf{(2)} saber qué contrato se firma en realidad (¿compraventa, obra o algo propio?); \textbf{(3)} cuando el edificio exista, ver cómo se describe el departamento (la \textit{unidad funcional}) en la escritura; y \textbf{(4)} si en lugar de un departamento se compra un terreno, cómo se describe el \textit{lote}.
\end{example}

% ------------------------------------------------------------
\chapter{La venta de cosas futuras}
\markboth{La venta de cosas futuras}{La venta de cosas futuras}

\section{Unidad 1 · Apertura y plan de la clase}

La compraventa se está cerrando en lo relativo a la \textbf{cosa}, y la clase siguiente se destinará al \textbf{precio}. La clase comienza con el tema de la \textbf{venta de cosas futuras}.

\section{Unidad 2 · La venta de cosas futuras y la venta en pozo}

\subsection{2.1 · Qué es una cosa futura}

\begin{definition}{Cosa futura}
\textit{<<Cuando hablamos de una cosa futura, estamos hablando de algo que el día de hoy no existe, no está.>>}
\end{definition}

\subsection{2.2 · Cómo se ve en la práctica inmobiliaria}

En la práctica inmobiliaria es habitual ver ventas de cosas que no existen.

\begin{example}{El cartel gigante}
\textit{<<No vemos el edificio: vemos el terreno baldío, en general, y un cartel gigante que dice ``Próximamente\ldots{} unidades de dos ambientes\ldots{} unidades de tres con balcón''. No se ve muy frecuentemente, ¿no? No es algo raro.>>}
\end{example}

\subsection{2.3 · El encuadre jurídico no es claro}

\textit{<<La estructura jurídica de estos negocios no está clara. El encuadre jurídico de estos negocios no es claro, no obstante la frecuencia con la que los vemos>>}, lo cual \textit{<<de por sí no llama la atención>>}. Se establece una analogía con otro tema ya visto: \textit{<<Lo mismo que con el boleto: es frecuente, no es algo raro, y sin embargo tenemos posturas antagónicas. Lo mismo pasa con las ventas en pozo.>>}

\subsection{2.4 · Qué se entiende por «venta en pozo»}

\begin{definition}{Venta en pozo}
El término <<venta en pozo>> no surge de la ley; es algo que se entiende en la cotidianidad. Se habla de venta en pozo cuando se trata de un inmueble totalmente pelado, de una casa que está por demolerse y en cuyo lugar van a estar las futuras unidades, o incluso de una obra en curso (por ejemplo: \textit{<<ya tengo dos pisos y faltan los otros ocho>>}).
\end{definition}

Se puede negociar en cualquiera de estos momentos, pero en general, cuando se habla de venta en pozo, se hace referencia a estos supuestos. Cuando el edificio ya está terminado y se entrega la posesión, es otro escenario, porque la cosa ya existe.

\begin{note}
\textbf{Delimitación del tema de la clase:} \textit{<<Hoy nos vamos a posicionar en el momento en donde la cosa todavía no existe, lo que me están ofreciendo.>>}
\end{note}

\section{Unidad 3 · ¿Por qué se dan estos negocios?}

La primera pregunta es: \textit{<<¿Por qué se dan estos negocios, o por qué son habituales?>>} Frente a la duda de por qué alguien querría comprar un terreno en lugar de ver la unidad terminada, se explica que estos negocios se dan básicamente por dos cuestiones.

\subsection{Ventaja del comprador: un menor precio}

El que compra seguramente paga menos dinero que lo que pagaría con una unidad terminada. \textit{<<Para esta persona, es un buen negocio desde una visión económica. Ahora, asume el riesgo de que la cosa no exista. Y por eso es el menor valor: en función del riesgo.>>}

\subsection{Ventaja del vendedor: financiamiento}

¿Por qué no espera el constructor a terminar el edificio y evita lidiar durante años con los compradores? Porque de esa forma puede financiar el proyecto: utiliza esos fondos para financiarse y, de lo contrario, no puede construir la edificación.

\begin{important}
\textbf{Conclusión:} \textit{<<No hay duda de que la cosa no existe y de que estos negocios se dan. Las ventajas podrían ser estas; puede haber otras, pero en principio de estas tres cosas no tenemos duda.>>}

Las tres ideas: (1) el comprador paga menos; (2) el comprador asume el riesgo de que la cosa no exista; (3) el constructor se financia con el dinero del comprador.
\end{important}

\section{Unidad 4 · ¿Se puede vender una cosa que no existe?}

Pasando al encuadre jurídico, la primera pregunta es si se puede vender una cosa que no existe. La respuesta es que \textbf{sí}: jurídicamente es posible vender una cosa que no existe, y el Código plantea \textbf{dos escenarios posibles}.

\begin{formula}{Código Civil y Comercial, art.~1131 (cosa futura)}
\textit{<<Si se vende cosa futura, se entiende sujeta a la condición suspensiva de que la cosa llegue a existir.}

\textit{El vendedor debe realizar las tareas y esfuerzos que resulten del contrato o de las circunstancias para que ésta llegue a existir en las condiciones y tiempo convenidos.}

\textit{El comprador puede asumir, por cláusula expresa, el riesgo de que la cosa no llegue a existir sin culpa del vendedor.>>}
\end{formula}

\subsection{Primer escenario · Compraventa sujeta a condición suspensiva}

El primer escenario consiste en sujetar el contrato de compraventa a una condición suspensiva: que la cosa llegue efectivamente a existir. Corresponde recordar lo visto en Obligaciones sobre las condiciones suspensivas y resolutorias (CCyC, arts.~343 y ss.).

\begin{note}
\textbf{Repaso: condición suspensiva y condición resolutoria.}

\textbf{Suspensiva:} \textit{<<cuando yo todavía no estoy obligado a ejecutar mis obligaciones hasta tanto ocurra algo. Por eso suspensiva: dejamos en suspenso. Mi obligación está en suspenso.>>}

\textbf{Resolutoria:} \textit{<<cuando yo dejaba de estar obligado>>}; la obligación se resuelve, se pone fin a ella.
\end{note}

Se puede vender una cosa que hoy no existe, pero el contrato queda condicionado a que la cosa llegue a existir. Si la cosa llega a existir, ambas partes deben ejecutar las obligaciones que asumieron: uno, transmitir el dominio, y el otro, pagar el precio en dinero. Si la cosa no llega a existir, no se deben nada.

\subsection{Segundo escenario · El comprador asume el riesgo: contrato aleatorio}

La otra posibilidad es que el comprador asuma el riesgo de que la cosa no exista nunca, con lo cual está obligado desde hoy al pago del precio. Se está entonces ante un \textbf{contrato aleatorio}.

En el primer caso, el comprador no asume ningún riesgo: si la cosa existe, se cumple. En el otro supuesto, el comprador asume el riesgo de que la cosa nunca llegue a existir y está obligado al pago desde el momento inicial; por eso es un contrato aleatorio: \textit{<<hoy contrato, estoy obligado, pero no sé si voy a tener retribución>>}.

\begin{formula}{Contratos aleatorios (CCyC, art.~968)}
\textit{<<Los contratos a título oneroso son conmutativos cuando las ventajas para todos los contratantes son ciertas. Son aleatorios, cuando las ventajas o pérdidas, para uno de ellos o para todos, dependen de un acontecimiento incierto.>>}
\end{formula}

\begin{table}[H]
\centering
\caption{Los dos escenarios del art.~1131 del CCyC}
\small
\begin{tabularx}{\textwidth}{>{\RaggedRight\arraybackslash}p{2.8cm} >{\RaggedRight\arraybackslash}X >{\RaggedRight\arraybackslash}X}
\toprule
 & \textbf{Condición suspensiva} & \textbf{Asunción del riesgo (aleatorio)} \\
\midrule
\textbf{Fundamento} & Art.~1131, 1.er párr.: regla general para la cosa futura. & Art.~1131, último párr.: debe pactarse por cláusula expresa. \\
\addlinespace
\textbf{Si la cosa llega a existir} & Ambas partes cumplen: se transmite el dominio y se paga el precio. & Ambas partes cumplen: se transmite el dominio; el precio ya estaba obligado. \\
\addlinespace
\textbf{Si la cosa no llega a existir} & No se deben nada: las obligaciones quedaron en suspenso. & El comprador igualmente debe el precio: asumió el riesgo. \\
\addlinespace
\textbf{Riesgo del comprador} & Ninguno. & Total: obligado al pago desde el momento inicial. \\
\addlinespace
\textbf{Naturaleza del contrato} & Conmutativo sujeto a condición. & Aleatorio. \\
\bottomrule
\end{tabularx}
\end{table}

Eso es lo que dice el Código sobre las ventas de cosas futuras: \textit{<<Esa es la teoría.>>}

\section{Unidad 5 · Aplicación a los inmuebles: ¿encaja alguno de los dos escenarios?}

Corresponde ver cómo se aplica lo anterior a los inmuebles.

\begin{example}{El departamento del cuarto piso}
\textit{<<Pensemos la situación: un terreno baldío. Ella [la desarrolladora] me vende a mí la futura unidad: el cuarto piso, que va a dar al frente, que va a tener balcón. ¿Cómo es el contrato?>>}
\end{example}

El razonamiento es el siguiente: el Código plantea, por un lado, la condición suspensiva (el comprador no está obligado a dar nada) y, por otro, el contrato aleatorio (el comprador asume el riesgo de que la desarrolladora no termine de construir). \textit{<<Pareciera que ninguno de los dos se ajusta.>>} Ni un caso ni el otro; la cuestión queda en suspenso para seguir pensando.

\begin{important}
\textbf{Por qué no encaja.} El vendedor \textbf{pide el dinero ahora}, por lo que no hay condición suspensiva pura (el comprador ya está pagando); y tampoco es un aleatorio puro, porque no puede decirse que, si el vendedor no construye, el comprador no pueda reclamarle nada.
\end{important}

Hay distintas posturas sobre esto: hay quienes consideran que es un contrato de compraventa, hay quienes consideran que se trata de un contrato de obra y hay otros que consideran que se trata de un contrato innominado. Ya se vio lo que dice la letra de la ley respecto a las ventas de cosas futuras en general; a continuación se ve lo que es un contrato de obra.

% ------------------------------------------------------------
\chapter{Naturaleza jurídica de la venta en pozo}
\markboth{Naturaleza jurídica de la venta en pozo}{Naturaleza jurídica de la venta en pozo}

\section{Unidad 6 · Compraventa y contrato de obra: dar y hacer}

El contrato de obra es aquel en que una parte se obliga a desarrollar determinada construcción o determinada tarea a cambio de un precio. La obligación principal en el contrato de obra es el \textbf{hacer}: la ejecución de la obra.

\begin{formula}{Contrato de obra (CCyC, art.~1251)}
\textit{<<Hay contrato de obra o de servicios cuando una persona, según el caso el contratista o el prestador de servicios, actuando independientemente, se obliga a favor de otra, llamada comitente, a realizar una obra material o intelectual o a proveer un servicio mediante una retribución.>>}
\end{formula}

\begin{note}
La retribución es elemento de la definición legal; por eso el contrato de obra es oneroso. (Nota editorial incluida en las notas originales.)
\end{note}

La definición que da el Código para el contrato de obra es bastante clara, al igual que las definiciones que da respecto a las ventas de cosas futuras.

\subsection{La diferencia sustancial}

Mientras que en el contrato de obra la obligación principal es el \textbf{hacer} (también podría decirse el dar, pero la obligación de dar o de entregar es solo la consecuencia natural de haber terminado, y lo primero es el hacer), en el contrato de compraventa la obligación principal es el \textbf{dar}: entregar la cosa y transferir el dominio.

\begin{table}[H]
\centering
\caption{Compraventa y contrato de obra}
\small
\begin{tabularx}{\textwidth}{>{\RaggedRight\arraybackslash}p{3.6cm} >{\RaggedRight\arraybackslash}X >{\RaggedRight\arraybackslash}X}
\toprule
 & \textbf{Contrato de compraventa} & \textbf{Contrato de obra} \\
\midrule
\textbf{Obligación principal} & \textsc{Dar}: transmitir el dominio, entregar la cosa. & \textsc{Hacer}: ejecutar la obra. \\
\addlinespace
\textbf{Entrega de la cosa} & Es la prestación central. & Es solo la consecuencia natural de haber terminado. \\
\addlinespace
\textbf{Lo que interesa al adquirente / comitente} & Recibir la cosa. & El modo en que se hace la cosa. \\
\bottomrule
\end{tabularx}
\end{table}

Queda abierta la pregunta: como adquirente o compradora, ¿interesa más recibir una cosa o el hacer de ella? \textit{<<En realidad van de la mano, porque el hacer también es hacer lo que queramos que haga. Si solo te interesa que te den la cosa, capaz te dan algo que no pactaste.>>}

\section{Unidad 7 · Primera postura: es un contrato de compraventa}

\textbf{Fundamento.} Quienes sostienen esta postura afirman que al adquirente solo le interesa que le entreguen el departamento futuro: no le importa qué proveedores contrate la desarrolladora, ni las cerámicas que ponga o deje de poner. La postura se centra en la entrega y, por ende, se trata de un contrato de compraventa.

\subsection{Críticas que recibe}

El Código plantea la situación de condición suspensiva o contrato aleatorio, y ninguno de los dos aplica, porque el vendedor pide el dinero ahora, y tampoco puede decirse que si no construye el comprador no pueda reclamarle nada. No se ajusta a ninguno de los dos escenarios.

\subsection{Cómo responden los defensores de la postura}

Responden que \textbf{no son normas de orden público}, y que en los contratos la regla general es que reina la \textbf{autonomía de la voluntad}: puede plantearse este contrato de compraventa con otras modalidades.

\subsection{Por qué, según esta postura, no es contrato de obra}

Porque el adquirente no tiene injerencia en la toma de decisiones de la desarrolladora: no puede decirle <<no contrate esa madera, que es malísima>>. La dirección de obra y todo eso lo encarga y lo ejecuta ella; el adquirente no tiene injerencia en el hacer.

Frente a la pregunta de si esta postura convence, la respuesta del curso fue negativa, y la clase continúa con la siguiente postura.

\section{Unidad 8 · Segunda postura: es un contrato de obra}

Esta postura sostiene lo que los propios alumnos adelantaron: difícilmente pueda afirmarse que al adquirente no le importa el hacer.

\textbf{Fundamento.} \textit{<<No me da lo mismo que me pongan el baño en el medio del living a que me lo pongan en otro lado. Y tan así que no nos da lo mismo a los adquirentes, que cada vez que suscribimos estos contratos se firman} \textbf{\textit{memorias descriptivas, planos, y se hace un detalle de los materiales a utilizar}}\textit{.>>}

Evidentemente, la ejecución no da lo mismo. Si bien el adquirente no va a estar llamando a la desarrolladora para saber si contrató a tal o cual persona, al suscribir el contrato suscribió el cronograma de obra, la memoria descriptiva, la calidad de los materiales, etcétera, evidenciando de esta manera que \textbf{sí le interesa el hacer}.

También en pequeños detalles (la cocina, por ejemplo), si el adquirente quiere poner una cosa, la pide: el adquirente tiene injerencia, puede aportar y decir <<yo quiero esto>>. Eso apoya la segunda postura.

\section{Unidad 9 · Tercera postura: contrato innominado y complejo}

Esta última postura sostiene que no es ni contrato de compraventa ni contrato de obra, sino un \textbf{contrato innominado y complejo}, que no se puede regir exclusivamente por las normas de uno ni de otro, porque sería forzar una situación jurídica a un encuadre jurídico que no le corresponde. Por lo tanto, todos los derechos y todas las obligaciones corren de acuerdo al contrato suscripto.

% TODO: contenido incompleto o ambiguo en las notas originales (un alumno formula una consulta en este punto; el audio no la registra)

\subsection{Cuadro comparativo de las tres posturas}

\begin{table}[H]
\centering
\caption{Las tres posturas sobre la venta en pozo}
\small
\begin{tabularx}{\textwidth}{>{\RaggedRight\arraybackslash}p{2.6cm} >{\RaggedRight\arraybackslash}X >{\RaggedRight\arraybackslash}X}
\toprule
\textbf{Postura} & \textbf{Fundamento} & \textbf{Crítica / respuesta} \\
\midrule
\textbf{1 · Compraventa} & Al adquirente solo le interesa recibir el departamento (dar); no le importan proveedores ni materiales. & \textit{Crítica:} ni la condición suspensiva ni el aleatorio del art.~1131 se ajustan (el vendedor cobra ahora).\par \textit{Respuesta:} no son normas de orden público; rige la autonomía de la voluntad. \\
\addlinespace
\textbf{2 · Contrato de obra} & Al adquirente le interesa el hacer: se firman memorias descriptivas, planos, cronograma y detalle de materiales; puede pedir pequeños cambios. & \textit{Crítica:} el adquirente no tiene injerencia en la dirección de obra ni en las decisiones del constructor. \\
\addlinespace
\textbf{3 · Innominado y complejo} & No puede regirse solo por las normas de la compraventa ni de la obra: sería forzar el encuadre jurídico. & Los derechos y obligaciones se rigen por lo pactado en el contrato suscripto. \\
\bottomrule
\end{tabularx}
\end{table}

% ------------------------------------------------------------
\chapter{Descripción del objeto: la unidad funcional}
\markboth{Descripción del objeto: la unidad funcional}{Descripción del objeto: la unidad funcional}

\section{Unidad 10 · Retomando la descripción del objeto}

Se retoma la descripción del objeto, tema arrastrado de la clase anterior, en la que se había comenzado con la unidad funcional y se había empezado a determinar qué datos interesan y cuáles se descartan.

\begin{note}
\textbf{Repaso: datos que sí hacen a la identificación de la unidad funcional.}
\begin{itemize}
    \item Calle, número y entre calles.
    \item Número de la unidad.
    \item Piso.
    \item Porcentual.
    \item Nomenclatura catastral.
\end{itemize}
Todos los demás datos no hacen a la identificación.
\end{note}

\begin{important}
\textbf{Objetivo del ejercicio:} elaborar la \textbf{cláusula primera de todo contrato de compraventa, que es la identificación del objeto}: A le vende a B tal cosa.
\end{important}

\section{Unidad 11 · Material de estudio}

Los materiales de estudio de la cátedra son \textbf{sugeridos}: pueden utilizarse o no. Hay archivos que tratan sobre el mismo tema, pero corresponden a dos autores distintos. Se enviará además un modelo de escritura para ver cómo se redacta la parte del inmueble.

\section{Unidad 12 · Redacción de la unidad funcional}

\begin{exercise}{Datos de la unidad funcional}
\textbf{Unidad:} cinco. \textbf{Piso:} primero. \textbf{Superficie total:} 50~\mcuad. \textbf{Porcentual:} 10\,\%. \textbf{Calle:} Av.\ Figueroa Alcorta n.º~333, entre Udaondo y Av.\ del Libertador. \textbf{Nomenclatura catastral:} circunscripción 15, sección 16, manzana 3, parcela 9.
\end{exercise}

Existen dos tipos de inmuebles: \textbf{unidad funcional o lote}. No importa si el lote está baldío, construido, etcétera: en un lote se escribe el perímetro.

\subsection{El intento de los alumnos}

\begin{note}
Redacción propuesta por un alumno: \textit{<<A le vende a B, bajo el régimen de propiedad horizontal, la unidad funcional número cinco, ubicada en el primer piso del inmueble de esta ciudad, con frente a la calle\ldots{} La unidad de la presente cuenta con una superficie cubierta y total de\ldots{} \mcuad, correspondiendo un porcentual de dominio del 10\,\%.>>} Valoración del docente: \textit{<<Redondo. Eso sí me encanta.>>}
\end{note}

\subsection{De dónde se toman los datos}

Cuando se escribe una unidad funcional por primera vez, se parte siempre del \textbf{reglamento de propiedad horizontal}. En el reglamento hay un plano con una grilla que numera todas las unidades del edificio, consigna las superficies parciales y totales y el porcentual que le corresponde. Al hacer una descripción por primera vez, hay que volcar todos estos datos al papel.

\begin{example}{Desglose de la superficie}
Se supone que la unidad tiene 50~\mcuad{} de total: 45~\mcuad{} cubiertos y 5~\mcuad{} descubiertos, y un porcentual del 10. En la planilla de superficies probablemente se tenga el desglose de cubierta, semicubierta y total. El total tiene que dar la sumatoria de todo. Hay veces que no da, y no es problema: hay que rectificar, etcétera.
\end{example}

\subsection{Criterio de orden: de lo más chico a lo más grande}

\begin{important}
\textbf{Regla práctica.} Al escribir una unidad funcional se parte \textbf{de lo más chiquito a lo más grande}. Es una cuestión de orden y de estrategia para que se entienda bien lo que se hace; casi todos los escribanos redactan de la misma manera, por una cuestión de lógica.
\end{important}

En la descripción habitual, se parte del número de la unidad funcional, que es lo más chiquito dentro del edificio: unidad cinco, que está en el primer piso (en el primer piso puede haber tres unidades o más). Después, el edificio en esta ciudad, que es más grande; después, entre calle tal y tal. Siguiendo siempre el mismo orden y el mismo criterio lógico, no se olvida ningún dato. Luego se pasa a lo particular de la unidad funcional, que es básicamente superficie y porcentual.

\subsection{El orden no es obligatorio, pero conviene mantenerlo}

Si bien este es el orden habitual, no es necesario que se redacte de esta manera: mientras los datos identificatorios estén consignados, no hay inconveniente, y nadie puede decir que está mal redactado.

Sin embargo, por una cuestión de \textbf{prolijidad}, en el sentido de evitar olvidos, conviene poner todo en el mismo orden, porque así se detecta más fácilmente en la lectura que falta algo. Se está acostumbrado a leer la escritura de cierta manera: se advierte enseguida que falta, por ejemplo, la comparecencia o el estado civil. Si se pone todo en el mismo orden, también con los inmuebles, no falta nada o, si falta, es fácil darse cuenta.

\subsection{Modelo de redacción dictado en clase}

\begin{example}{Unidad funcional (versión dictada)}
\textit{<<A transmite a título de venta a B>> (o <<A vende a B>>) <<el inmueble descripto de la siguiente manera: la unidad funcional cinco, ubicada en el primer piso, integrante del edificio sito en esta ciudad con frente a la avenida Figueroa Alcorta, número 333, entre Udaondo y avenida del Libertador. Nomenclatura catastral: circunscripción 15, sección 16, manzana 3, parcela 9. La unidad consta de una superficie cubierta de 45~\mcuad, una superficie descubierta de 5~\mcuad, lo que hace una superficie total de 50~\mcuad, y le corresponde un porcentual de [10\,\%].>>}
\end{example}

\textbf{Aclaraciones.} Si es una cochera (si es una unidad funcional) es lo mismo. Si es una unidad complementaria, es la misma descripción; solo que, en vez de consignar <<unidad funcional>>, se consigna <<unidad complementaria>>.

No hay que marearse con los datos y el orden: mientras estén todos los datos y se pongan en el orden elegido, está bien. Lo que siga en la lógica de quien redacta, está bien.

\section{Unidad 13 · ¿Números o letras? Art.~303 del CCyC y art.~61 de la Ley 404}

Cerrada la redacción, surge la pregunta: ¿qué va en letras? En el debate se plantearon posiciones diversas: que el número de la unidad funcional, el piso y la superficie total van en letras; que algunos consideran que todo lo que forme parte de la superficie (cubierta, descubierta, total) va en letras y otros que solo la total; que en escrituras muy viejas o de provincia se ha visto el número de la calle con letras, y que la nomenclatura y el porcentual se han visto en números. Sobre la unidad funcional y la superficie total \textit{<<nadie dice otra cosa>>}; lo que diga el reglamento.

\textbf{¿Por qué se hace esto?} Porque no puede ponerse como se le ocurra: hay \textbf{normativa que manda a poner en letras ciertos números}. Surge de un artículo del Código y de un artículo de la Ley 404 (ley orgánica notarial de la Ciudad de Buenos Aires); es un artículo cada uno.

\subsection{El artículo del Código Civil y Comercial}

\begin{formula}{Código Civil y Comercial, art.~303 (uso de números)}
\textit{<<Pueden usarse números, excepto para las cantidades que se entregan en presencia del notario y para las que corresponden a elementos esenciales del acto jurídico.>>}
\end{formula}

\subsection{El artículo de la Ley 404}

\begin{formula}{Ley 404 (CABA), art.~61 (requisitos generales de los documentos notariales)}
\textit{<<Todos los documentos deberán ser escritos sin espacios en blanco en su texto. No se emplearán abreviaturas ni iniciales, excepto cuando consten en los documentos que se transcriben, se trate de constancias de otros documentos o de signos o abreviaturas científicas o socialmente admitidas con sentido unívoco. No se utilizarán guarismos para expresar el número de escritura, su fecha, el precio o monto de la operación, las cantidades entregadas en presencia del escribano ni las condiciones de pago.>>}
\end{formula}

% TODO: contenido incompleto o ambiguo en las notas originales (el audio capta de manera confusa el texto del art. 61; las notas originales indican que el texto reproduce el artículo con la corrección correspondiente)

Por eso el número de escritura y la fecha siempre van en letras, \textit{<<sí o sí>>}. El precio o monto de la operación se trata en la clase siguiente, y las cantidades entregadas en presencia o las condiciones de pago, después.

\subsection{Consecuencia para la unidad funcional}

Los elementos esenciales del acto jurídico no pueden describirse en números: la unidad funcional es justamente el elemento esencial del acto jurídico, por lo que no hay duda de que va en letras. La superficie total también va en letras. El porcentual de dominio tampoco admite duda: por eso se ve siempre <<diez enteros>>, <<diez enteros con cincuenta centésimos>>.

\begin{table}[H]
\centering
\caption{Qué va en letras y qué en números}
\small
\begin{tabularx}{\textwidth}{>{\RaggedRight\arraybackslash}p{4cm} >{\RaggedRight\arraybackslash}p{2.9cm} >{\RaggedRight\arraybackslash}X}
\toprule
\textbf{Dato a consignar} & \textbf{Forma de escritura} & \textbf{Observación} \\
\midrule
\textbf{Unidad funcional (número) y piso} & En \textsc{letras} & <<No hay duda que va en letras>>: elemento esencial del acto. \\
\addlinespace
\textbf{Superficie total} & En \textsc{letras} & <<Lo importante es la total.>> \\
\addlinespace
\textbf{Porcentual de dominio} & En \textsc{letras} & Por ejemplo, <<diez enteros>>; <<diez enteros con cincuenta centésimos>>. \\
\addlinespace
\textbf{Superficies parciales (cubierta, descubierta)} & Libre: números o letras & <<Hay quienes la consignan en números, hay quienes en letras; la pueden consignar como quieran.>> \\
\addlinespace
\textbf{Nomenclatura catastral} & En \textsc{números} & <<Todo el mundo la pone en números.>> \\
\addlinespace
\textbf{Número de escritura y fecha} & En \textsc{letras} (art.~61, Ley 404) & Siempre, <<sí o sí>>. \\
\bottomrule
\end{tabularx}
\end{table}

\begin{important}
\textbf{Regla de oro para el examen.} Ante la duda de qué va en letras y qué en números, equivocarse está mal: si iba en letras y se puso en números, es un \textbf{error fuerte}, porque quien corrige lo considera un error. Ante la duda, conviene poner siempre en letras, porque de esa manera se zanja la cuestión: nadie dirá <<pero podrías haber puesto números>>. Para las superficies parciales, como se prefiera: en números o en letras.
\end{important}

\section{Unidad 14 · Cómo leer la consigna del examen}

Si se equivoca en esto (por ejemplo, si se pone la unidad funcional como <<5>>), en una escritura real, o en el examen para ser escribano, es un \textbf{error fuerte}. Puede pensarse que el resto de la escritura estaba bien, pero justamente por ser un \textbf{elemento esencial} el error es más grave que uno cometido en una cláusula común.

Al rendir el examen, la consigna aporta datos que en una primera lectura pueden parecer intrascendentes pero que \textbf{hacen a la diferencia en la nota final}, porque son detalles de atención.

\begin{example}{Detalles de la consigna}
Si la consigna dice que es un departamento con patio o con balcón, evidentemente tiene superficie descubierta. Eso puede no hacer aprobar o desaprobar, pero si se está entre un siete y un ocho, puede marcar la diferencia. Además, si se indica que es un departamento, obviamente no puede escribirse un lote.
\end{example}

Estas cosas ocurren porque uno está nervioso y se apura, y hay detalles que se escapan igual.

% ------------------------------------------------------------
\chapter{Descripción del objeto: el lote}
\markboth{Descripción del objeto: el lote}{Descripción del objeto: el lote}

\section{Unidad 15 · Del gráfico al papel: qué datos exige un lote}

Se parte de un gráfico (el plano dibujado en el pizarrón) que hay que pasar al papel. Un primer intento de los alumnos resultó confuso en cuanto a orientaciones y linderos.

\begin{note}
\textbf{Primer intento de los alumnos:} \textit{<<Un inmueble en esta ciudad, con frente a la calle número 333, que se designa como\ldots{} mide por su frente al noroeste\ldots{} 8 metros lindando con la calle Figueroa Alcorta; por su contrafrente, 8 metros lindando con el lote 3; por su costado al este, 10 metros, con la calle Libertador; y por el oeste, su otro costado, 10 metros, lindando con la parcela catastral 15.>>} El docente señala la confusión y propone avanzar despacio.
\end{note}

Hay que \textbf{identificar la parcela} sobre la cual está el predio. La circunscripción, sección, manzana y parcela: la parcela es el lote; la manzana es la cuadra donde se está.

\begin{important}
\textbf{Datos que sí o sí se necesitan para escribir un lote:}
\begin{itemize}
    \item Número de lote (y manzana).
    \item Medidas y linderos.
    \item Orientación.
    \item Calle, número y entre calles.
    \item Nomenclatura catastral.
\end{itemize}
\end{important}

\section{Unidad 16 · Cómo se lee el plano del lote}

El plano del inmueble muestra, primero, la calle a la que da. Debajo de esa línea hay un número, que indica cuánto mide ese lado.

\begin{example}{El plano}
Aparecen circulitos que indican la designación de los lotes. El lote propio es el siete, y tiene al costado el cinco, abajo el tres y del otro lado el cuatro: son los \textbf{lotes linderos}. En cada lado figura la medida del lado: diez metros. Si un lado no tiene número, es porque mide lo mismo que el opuesto: este va a medir ocho también, y este va a medir diez.
\end{example}

Para describir un lote, además del número y de la nomenclatura, hay que consignar la medida de cada lado (el \textbf{perímetro}), la orientación a la que da cada lado (norte, sur, este, oeste) y qué lote es lindero. Es decir: \textbf{medidas de los lados, lotes linderos y orientación de cada lado}.

\subsection{El sentido de la descripción}

En general se arranca del \textbf{frente} al escribir y se sigue en el \textbf{sentido de las agujas del reloj}, hasta completar todo el círculo. No es necesario que sea así, y puede describirse en otro orden; pero por una cuestión de prolijidad se lleva en el sentido de las agujas del reloj.

\subsection{Las orientaciones}

Junto al plano hay una \textbf{flecha}, que siempre apunta al norte; por eso el inmueble del ejemplo tiene frente al norte.

\begin{note}
Nota editorial de las notas originales: el audio menciona a la avenida del Libertador como calle del frente; en la redacción dictada a continuación, el frente da a la avenida Figueroa Alcorta, que es la que se mantiene en este documento.
\end{note}

% TODO: contenido incompleto o ambiguo en las notas originales (discrepancia en el audio sobre la calle del frente)

\section{Unidad 17 · Redacción de la descripción del lote}

\begin{example}{Lote (versión dictada y ordenada)}
\textit{<<A transmite a B el inmueble sito en esta ciudad, con frente a la avenida Figueroa Alcorta, número trescientos treinta y tres, entre Udaondo y avenida del Libertador, que se designa como lote siete de la manzana nueve.}

\textit{El inmueble consta de las siguientes medidas y linderos: por su frente, al norte, linda con la avenida Figueroa Alcorta y mide ocho metros; por su costado al oeste, linda con la parcela cinco y mide diez metros; por su contrafrente, al sur, mide ocho metros y linda con el lote tres; y por su costado al este, mide diez metros y linda con el lote cuatro.>>}
\end{example}

Método: en cada lado se dicen \textbf{tres cosas: orientación, lindero y metros}. Por ejemplo: por su frente al norte, linda con la avenida y mide ocho metros. Aquí son metros, no metros cuadrados: son los \textbf{metros lineales}. Hasta ahí puede cortarse y decir <<el inmueble consta de las siguientes medidas y linderos>> o <<medidas y linderos>>, como se prefiera.

\section{Unidad 18 · ¿Por qué se nombra la calle solo en el frente?}

La calle se nombra en el primer lado (el del norte) y no se vuelve a nombrar en los otros porque el frente tiene que dar a la calle: no puede estar en el medio de la manzana, y sí o sí el frente linda con la calle. En los demás lados no se nombran calles porque, si tienen número de lote, es que no lindan con calles. Los otros costados pueden pasar a lindar con una calle si el lote está en una esquina: se consigna porque hay que decir lo que se ve en el plano, y en una esquina lo que se ve es una calle.

\begin{important}
\textbf{Regla práctica.} El \textbf{frente} siempre linda con la calle. Los demás lados lindan con \textbf{otros lotes} (se nombran por su número) salvo que el lote sea de \textbf{esquina}, caso en que un costado también linda con una calle.
\end{important}

\section{Unidad 19 · Superficie, nomenclatura y tarea para la clase siguiente}

Después de las medidas, la descripción consigna una superficie total de tantos metros cuadrados. En el lote \textbf{no hay porcentual}, porque el porcentual va solo para la propiedad horizontal. Las medidas de los lados son en metros, no en metros cuadrados, porque son lineales.

\subsection{La superficie total del lote}

Si la consigna dice que la casa tiene 50~\mcuad{} cubiertos, o que la descripción para la venta dice eso, no significa que el lote tenga 50~\mcuad, porque lo que interesa es el perímetro que se vende.

\begin{note}
Nota de estudio incluida en las notas originales: con el ejemplo dictado (frente de 8~m por costado de 10~m), la superficie total del lote sería de 80~\mcuad. El docente no calcula este valor en el audio.
\end{note}

\subsection{La nomenclatura catastral}

Finalmente, se consigna la nomenclatura catastral.

\subsection{Tarea para la clase siguiente}

\begin{exercise}{Descripción de una unidad funcional}
Se deja para que se haga una descripción, que se corrige la clase siguiente.

\textbf{Datos dictados:} unidad funcional uno, planta baja; superficie cubierta 50~\mcuad, semicubierta 5~\mcuad; porcentual <<con 75>> (centésimos); dirección sobre una calle <<entre>> dos calles (los nombres no resultan audibles en la grabación).
\end{exercise}

\begin{note}
Nota editorial de las notas originales: 50~\mcuad{} cubiertos más 5~\mcuad{} semicubiertos suman 55~\mcuad; el audio registra <<60>> como total, un dato inconsistente que conviene verificar con la consigna original. Como se explicó en la Unidad~12: <<el total tiene que dar la sumatoria; hay veces que no da, y no es problema: hay que rectificar>>.
\end{note}

% TODO: contenido incompleto o ambiguo en las notas originales (total de superficie y porcentual de la tarea)

% ============================================================
% CORNELL
% ============================================================
\chapter*{Resumen Cornell}
\addcontentsline{toc}{chapter}{Resumen Cornell}
\markboth{Resumen Cornell}{Resumen Cornell}

A la izquierda figuran las preguntas y conceptos clave para autoevaluarse; a la derecha, las respuestas; al pie, la síntesis integradora de la clase.

{\small
\begin{longtable}{>{\RaggedRight\arraybackslash}p{4.2cm} >{\RaggedRight\arraybackslash}p{10.3cm}}
\toprule
\textbf{Preguntas / palabras clave} & \textbf{Notas / respuestas} \\
\midrule
\endfirsthead
\toprule
\textbf{Preguntas / palabras clave} & \textbf{Notas / respuestas} \\
\midrule
\endhead
\midrule
\multicolumn{2}{r}{\textit{continúa en la página siguiente}} \\
\endfoot
\endlastfoot

\textbf{¿Qué es una cosa futura y qué es la venta en pozo?} &
Cosa futura es la que \textbf{hoy no existe}. La venta en pozo es la que se hace sobre un terreno baldío, una casa por demoler o una obra en curso (por ejemplo, dos pisos hechos y ocho por construir). Cuando el edificio ya está terminado y se entrega la posesión, la cosa ya existe y es otro escenario. \\
\addlinespace
\cmidrule{1-2}
\textbf{¿Por qué se realizan estos negocios?} &
\textbf{Comprador:} paga menos, a cambio de asumir el riesgo de que la cosa no exista. \textbf{Constructor:} se financia con el dinero de los compradores; sin esos fondos no podría construir. \\
\addlinespace
\cmidrule{1-2}
\textbf{¿Se puede vender una cosa que no existe?} &
\textbf{Sí.} El Código lo permite (art.~1131 CCyC) y plantea dos escenarios: condición suspensiva o asunción del riesgo. \\
\addlinespace
\cmidrule{1-2}
\textbf{¿Qué dice el art.~1131 CCyC?} &
La venta de cosa futura se entiende sujeta a la \textbf{condición suspensiva} de que la cosa llegue a existir. El vendedor debe hacer las tareas y esfuerzos necesarios para que exista en las condiciones y tiempo convenidos. El comprador puede asumir, por \textbf{cláusula expresa}, el riesgo de que no llegue a existir sin culpa del vendedor. \\
\addlinespace
\cmidrule{1-2}
\textbf{Condición suspensiva vs.\ contrato aleatorio} &
\textbf{Suspensiva:} si la cosa existe, las partes cumplen (dominio y precio); si no, no se deben nada. \textbf{Aleatorio:} el comprador queda obligado al pago desde el inicio, aunque la cosa nunca exista; las ventajas o pérdidas dependen de un acontecimiento incierto (art.~968). \\
\addlinespace
\cmidrule{1-2}
\textbf{¿Por qué no encajan en el departamento en pozo?} &
Porque el vendedor \textbf{cobra ahora} (no hay una suspensión pura) y tampoco puede decirse que, si no construye, el comprador no puede reclamarle nada (no es un aleatorio puro). \\
\addlinespace
\cmidrule{1-2}
\textbf{¿Cuál es la diferencia entre compraventa y contrato de obra?} &
En la compraventa la obligación principal es \textbf{dar} (transmitir el dominio). En la obra es \textbf{hacer} (ejecutar la obra); la entrega es la consecuencia natural de haber terminado (art.~1251 CCyC). \\
\addlinespace
\cmidrule{1-2}
\textbf{¿Cuáles son las tres posturas sobre la venta en pozo?} &
\textbf{(1) Compraventa:} al adquirente solo le interesa recibir la unidad; las normas del art.~1131 no son de orden público y rige la autonomía de la voluntad. \textbf{(2) Obra:} al adquirente sí le importa el hacer (memorias, planos, cronograma, materiales). \textbf{(3) Innominado y complejo:} no se rige solo por las normas de uno ni de otro; se aplica lo pactado. \\
\addlinespace
\cmidrule{1-2}
\textbf{¿Qué es la cláusula primera de la compraventa?} &
La \textbf{identificación del objeto}: <<A le vende a B tal cosa>>. El inmueble puede ser una \textbf{unidad funcional o un lote}. \\
\addlinespace
\cmidrule{1-2}
\textbf{¿Qué datos identifican una unidad funcional y en qué orden se redactan?} &
Número de unidad, piso, edificio, calle, número y entre calles, nomenclatura catastral, superficies y porcentual. Se parte \textbf{de lo más chico a lo más grande} y luego se vuelve a lo particular (superficie y porcentual). El orden no es obligatorio, pero conviene por prolijidad y para detectar olvidos. Los datos salen del \textbf{reglamento de propiedad horizontal}. \\
\addlinespace
\cmidrule{1-2}
\textbf{¿Qué va en letras y por qué?} &
Lo que es \textbf{elemento esencial del acto jurídico}: unidad funcional, piso, superficie total y porcentual de dominio (art.~303 CCyC). Número de escritura y fecha, también (art.~61 Ley 404). La nomenclatura catastral va en números. Ante la duda: \textbf{letras}. \\
\addlinespace
\cmidrule{1-2}
\textbf{¿Cómo se describe un lote?} &
Número de lote y manzana, calle, número y entre calles, nomenclatura catastral y \textbf{perímetro}: para cada lado, \textbf{orientación, lindero y medida}. Se empieza por el frente y se sigue en el sentido de las agujas del reloj. \\
\addlinespace
\cmidrule{1-2}
\textbf{¿Cómo se leen la orientación y los linderos?} &
La \textbf{flecha} del plano indica el norte. El \textbf{frente} siempre linda con la calle; los demás lados lindan con otros lotes (se nombran por su número), salvo en un lote de \textbf{esquina}, donde un costado también da a una calle. \\
\addlinespace
\cmidrule{1-2}
\textbf{¿Qué cuidados hay en el examen?} &
Leer los datos de la consigna (si hay patio o balcón, hay superficie descubierta; si es un departamento, no se describe un lote). En el lote \textbf{no hay porcentual} (solo propiedad horizontal), las medidas lineales van en \textbf{metros} y lo que interesa es el perímetro, no los metros cubiertos de la casa. \\
\midrule
\multicolumn{2}{p{14.5cm}}{\textbf{Resumen:} La clase muestra cómo el Derecho resuelve un negocio frecuente pero de encuadre discutido: la venta de inmuebles que todavía no existen. El art.~1131 del CCyC permite vender cosa futura bajo condición suspensiva o con asunción expresa del riesgo (contrato aleatorio), pero ninguna de las dos modalidades describe con exactitud la venta en pozo, donde el comprador paga hoy y el vendedor debe construir; de ahí las tres posturas (compraventa, contrato de obra o contrato innominado y complejo), que giran sobre si lo central es el \textit{dar} o el \textit{hacer}. Luego la clase pasa a la práctica notarial: la cláusula primera de la compraventa identifica el objeto, que puede ser una unidad funcional, redactada de lo más chico a lo más grande y con los elementos esenciales en letras (art.~303 CCyC y art.~61 Ley 404), o un lote, descrito por su perímetro con orientación, lindero y medida de cada lado. El hilo conductor es la precisión: saber qué se vende, cómo se identifica y qué se está firmando.} \\
\bottomrule
\end{longtable}
}

\end{document}
